\documentclass[11pt]{article}
\usepackage[a4paper,margin=1in]{geometry}
\usepackage{array}
\usepackage{booktabs}
\usepackage{enumitem}
\usepackage{fontspec}
\usepackage{tabularx}
\setmainfont{Charis SIL}
\newfontfamily\hanfont{BabelStone Han}
\newcommand{\han}[1]{{\hanfont #1}}
\setlength{\parindent}{0pt}
\setlength{\parskip}{0.6em}
\begin{document}
\section*{Pilot a/b subseries detection}
This pilot infers handwritten subgroup `a/b' labels from the Middle Chinese forms attached to each subgroup in the hand-written \texttt{main.tex}. It uses the subgroup head's own MC forms as the primary signal.
\begin{enumerate}[nosep]
\item Treat a subgroup as type `b' when its extracted MC forms show an `i'-medial after the onset, for example \texttt{gie}, \texttt{phiu}, or \texttt{ṅiəH}.
\item Treat a subgroup as type `a' when the extracted forms lack that `i'-medial, for example \texttt{ka} or \texttt{puH}.
\item Leave bare zero-onset \texttt{i-} forms unresolved; this ambiguity does not arise in the selected pilot entries.
\end{enumerate}
Pilot entries reviewed: \texttt{18-01}, \texttt{01-67}.
Matched handwritten label in 8/8 inspected subgroups.

\subsection*{18-01 (\han{可})}
Source lines 385--512.
{\small
\begin{tabularx}{\textwidth}{@{}>{\raggedright\arraybackslash}p{1.4cm}cc>{\raggedright\arraybackslash}X>{\raggedright\arraybackslash}X@{}}
\toprule
Head & Hand & Pred & MC forms & Handwritten RHS LaTeX \\
\midrule
\han{何} & \texttt{a} & \texttt{a} & ḫa, ḫaX & g\textbackslash\{\}textoverset\{a\}\{a\}y \\
\han{哥} & \texttt{a} & \texttt{a} & ka & k\textbackslash\{\}textoverset\{a\}\{a\}y \\
\han{奇} & \texttt{b} & \texttt{b} & gie, kie & q\textbackslash\{\}textoverset\{b\}\{a\}y \\
\bottomrule
\end{tabularx}
}
\begin{itemize}[nosep]
\item \han{何}: all extracted forms lack an i-medial after the onset.
\item \han{哥}: all extracted forms lack an i-medial after the onset.
\item \han{奇}: all extracted forms show an i-medial or a dedicated palatal onset.
\end{itemize}

\subsection*{01-67 (\han{父})}
Source lines 3036--3159.
{\small
\begin{tabularx}{\textwidth}{@{}>{\raggedright\arraybackslash}p{1.4cm}cc>{\raggedright\arraybackslash}X>{\raggedright\arraybackslash}X@{}}
\toprule
Head & Hand & Pred & MC forms & Handwritten RHS LaTeX \\
\midrule
\han{布} & \texttt{a} & \texttt{a} & puH & p\textbackslash\{\}textoverset\{a\}\{a\} \\
\han{浦} & \texttt{a} & \texttt{a} & phuX & p\textbackslash\{\}textoverset\{a\}\{a\}₂ \\
\han{捕} & \texttt{a} & \texttt{a} & buH & p\textbackslash\{\}textoverset\{a\}\{a\}₃ \\
\han{尃} & \texttt{b} & \texttt{b} & phiu & p\textbackslash\{\}textoverset\{b\}\{a\} \\
\han{旉} & \texttt{b} & \texttt{b} & phiu & p\textbackslash\{\}textoverset\{b\}\{a\}₂ \\
\bottomrule
\end{tabularx}
}
\begin{itemize}[nosep]
\item \han{布}: all extracted forms lack an i-medial after the onset.
\item \han{浦}: all extracted forms lack an i-medial after the onset.
\item \han{捕}: all extracted forms lack an i-medial after the onset.
\item \han{尃}: all extracted forms show an i-medial or a dedicated palatal onset.
\item \han{旉}: all extracted forms show an i-medial or a dedicated palatal onset.
\end{itemize}

\end{document}
